\chapter{Plugins}
\label{kap:format-plugins}

Ein Plugin ist eine Bibliothek (DLL, Shared Object oder Dylib) mit dem
Verhalten eines Bauteils, das \esim nicht selbst rechnet: ein Feldlöser für
Magnete und Motoren, ein Halbleitersimulator auf einem Gitter, ein
thermisches 3D-Modell, ein neuronales Netz. \esim modelliert, was um das
Bauteil herum geschieht (Ansteuerung, Verstärker, Regler, Mechanik); das
Plugin liefert die Gleichungen des Bauteils. Beide werden gemeinsam gelöst,
nicht nacheinander: Newton-Verfahren, Schrittweitenregelung,
Wechselstromanalyse und Rauschanalyse arbeiten über Schaltung und Plugin
hinweg.

\section{Aufbau eines Plugin-Bauteils}
\label{sec:plugin-aufbau}

Aus Sicht von \esim ist ein Plugin-Bauteil ein Bauteil mit fünf Arten von
Größen:

\begin{description}[style=nextline,leftmargin=2em]
  \item[Klemmen (\code{nodes})] Knoten der Schaltung. Das Plugin sieht ihre
    Spannungen und liefert den Strom (und die Ladung), der in das Bauteil
    fließt, wie ein SPICE-Bauteil. Ein Halbleiter hat hier seine Kontakte.
  \item[Eingänge (\code{inputs})] Ausdrücke der Eingabedatei, die das Plugin
    nur liest: eine Position, ein Winkel, eine Temperatur, eine
    Lichtleistung. Jeder Eingang des Plugins muss belegt werden.
  \item[Zweige] Ströme zwischen zwei Klemmen, etwa einer Wicklung. \esim
    legt dafür eine Unbekannte an und bucht den Strom in die Klemmenzeilen;
    die Zweiggleichung steht im Plugin.
  \item[Ausgänge (\code{outputs})] Größen, die das Plugin berechnet, zum
    Beispiel Kraft, Drehmoment oder Verlustleistung. Jeder Ausgang wird eine
    Variable der Eingabedatei und darf in Gleichungen, Verhaltensquellen,
    Ausgaben und Messgrößen stehen.
  \item[innere Unbekannte] Zustände des Plugins (Feldgrößen, Potentiale und
    Ladungsträgerdichten je Gitterpunkt). Sie sind Unbekannte des
    Gesamtsystems und werden mit der Schaltung gelöst.
\end{description}

Die Gleichungen des Plugins haben dieselbe Form wie die der Schaltung,
$\mathrm{d}q(x)/\mathrm{d}t + f(x,t) = 0$. Das Plugin liefert zu jeder
Auswertung $f$, $q$ und die Jacobi-Matrizen $\partial f/\partial x$ und
$\partial q/\partial x$.

\section{Einbinden}
\label{sec:plugin-einbinden}

Das Beispiel ist ein Hubmagnet. Das Plugin enthält Wicklung und Feld, die
Mechanik steht in der Eingabedatei. Die Position \code{z} geht als Eingang
ins Plugin, die Kraft kommt als Ausgang zurück:

\begin{json}
{
  "parameters": { "Rcoil": 2.0, "mass": 0.05, "k": 200, "c": 1.5, "zrest": "10m" },
  "plugins": {
    "magnet": { "library": "esim_mock_plugin",
                "config": { "kind": "solenoid", "L0": "10m", "z0": "5m" } }
  },
  "components": [
    { "type": "V", "name": "VS", "nodes": ["a", "0"],
      "wave": { "type": "pulse", "v1": 0, "v2": 5, "delay": "1m", "rise": "10u" } },
    { "type": "plugin", "name": "X1", "plugin": "magnet", "nodes": ["a", "0"],
      "config": { "R": "Rcoil" },
      "inputs":  { "z": "z" },
      "outputs": { "F": "Fm" } }
  ],
  "variables": { "z": { "init": "10m" }, "vz": { "init": 0 } },
  "equations": [
    "der(z) = vz",
    "mass*der(vz) = Fm - k*(z - zrest) - c*vz"
  ],
  "analyses": [ { "type": "tran", "tstop": "60m" } ],
  "outputs": ["z", "i(X1)", "Fm"]
}
\end{json}

\begin{itemize}
  \item \code{plugins} vereinbart Plugins mit Namen. \code{library} ist der
    Dateiname mit oder ohne Endung (\datei{.dll}, \datei{.so},
    \datei{.dylib}); relative Namen gelten zur Eingabedatei. \code{config}
    ist die Konfiguration, die das Plugin beim Anlegen bekommt (Modelldatei,
    Parameter, Gitter). Das Plugin legt fest, welche Schlüssel es kennt.
  \item Das bauteilweise \code{config} ergänzt die Konfiguration des Plugins
    und überschreibt gleichnamige Einträge. Zeichenketten, die eine Zahl
    oder ein konstanter Ausdruck der Parameter sind (\code{"4.7k"},
    \code{"Rcoil"}), werden zu Zahlen; alle anderen, etwa Dateinamen, bleiben
    Text. Damit laufen Parameterstudien und Monte-Carlo-Läufe auch durch
    Plugin-Parameter.
  \item Statt \code{"plugin": "magnet"} darf \code{"library"} direkt am Bauteil
    stehen; dann entfällt der Eintrag in \code{plugins}.
\end{itemize}

\section{Schlüssel des Bauteils}
\label{sec:plugin-schluessel}

\begin{longtable}{@{}P{2.8cm}P{3.6cm}P{7.4cm}@{}}
  \toprule
  Schlüssel & Wert & Bedeutung \\
  \midrule
  \endhead
  \bottomrule
  \endfoot
  \code{type} & \code{"plugin"} & Bauteiltyp \\
  \code{nodes} & Liste von Knotennamen & die Klemmen in der Reihenfolge des Plugins; die Anzahl muss stimmen \\
  \code{plugin} & Name aus \code{plugins} & oder \code{library} \\
  \code{library} & Dateiname & Bibliothek direkt am Bauteil \\
  \code{config} & Objekt & Konfiguration; ergänzt die des Plugins \\
  \code{inputs} & \code{{"eingang": "ausdruck"}} & Ausdruck je Eingang: Spannungen, Ströme, Variablen, Zeit, Parameter \\
  \code{outputs} & \code{{"ausgang": "variable"}} & Name der Variable je Ausgang; ohne Eintrag \code{<bauteil>.<ausgang>} \\
  \code{snapshots} & \code{true} | \code{{"times": [...]}} & Feldausgabe, Abschnitt~\ref{sec:plugin-schnappschuesse} \\
\end{longtable}

\esim bildet die Ableitungen der Eingangsausdrücke selbst
(Kettenregel). Das Plugin muss also nur die Ableitung nach seinen Eingängen
liefern, nicht nach den Größen der Schaltung.

\subsection*{Namen von außen}

\begin{longtable}{@{}P{4.2cm}P{9.6cm}@{}}
  \toprule
  Ausdruck & Bedeutung \\
  \midrule
  \endhead
  \bottomrule
  \endfoot
  \code{i(X1)} & Strom des ersten Zweigs von \code{X1} \\
  \code{i(X1.coil)} & Strom des Zweigs \code{coil}; der Name kommt vom Plugin \\
  \code{Fm} & Ausgang, wie unter \code{outputs} benannt \\
  \code{X1.F} & Ausgang ohne eigenen Namen in \code{outputs} \\
  \code{X1.y} & innere Unbekannte \code{y}; sie steht nicht in der Standardausgabe, wohl aber in \code{outputs} der Datei \\
\end{longtable}

\section{Verhalten in den Analysen}
\label{sec:plugin-analysen}

\begin{description}[style=nextline,leftmargin=2em]
  \item[\code{op}, \code{dc}] Das Plugin rechnet im Gleichgewicht. Startwerte
    der inneren Unbekannten kommen vom Plugin (zum Beispiel eine
    Gleichgewichtslösung). Die Hilfsverfahren zur Arbeitspunktsuche
    (gmin-Stepping, Quellenhomotopie, Pseudo-Transiente) wirken auf
    Unbekannte, die das Plugin als Spannung meldet.
  \item[\code{tran}] Die Zeitableitung bildet \esim aus der Ladung $q$ des
    Plugins. Das Plugin darf Zeitpunkte vorgeben, auf denen \esim landet,
    und die Schrittweite begrenzen, etwa auf einen Bruchteil einer
    Wirbelstrom-Zeitkonstante. Der Fehlerschätzer des Verfahrens arbeitet mit
    den Toleranzen der Unbekannten; Unbekannte in Volt (Potentiale) skalieren
    besser als Dichten.
  \item[\code{ac}] Die Kleinsignalantwort folgt aus $G+\mathrm{j}\omega C$ des
    Plugins am Arbeitspunkt. Ein Eingang, der von einer Spannung der Schaltung
    abhängt, wirkt über die Kettenregel mit. Licht als Eingang ist so mit einer
    Quelle mit \code{"ac": 1} an einem Knoten messbar. Ein
    frequenzabhängiges Verhalten im Plugin selbst ist (noch) nicht möglich.
  \item[\code{noise}] Meldet das Plugin Rauschquellen (Stromquellen zwischen
    zwei seiner Zeilen mit weißem und Funkelrauschen), rechnet \esim sie
    zusammen mit denen der Schaltung. Der Aufwand je Frequenz ist ein
    Gleichungssystem plus ein Skalarprodukt je Quelle.
  \item[\code{sens}, \code{param\_sweep}, \code{monte\_carlo}] Die Eingabedatei
    wird je Lauf neu gelesen, das Plugin also neu angelegt. Ein Plugin mit
    großem Gitter hält dieses in der Bibliothek, damit das nicht jedes Mal
    Zeit kostet. Ableitungen (\code{sens}) sind Differenzenquotienten.
  \item[Newton-Verfahren] Das Plugin darf wie SPICE bei Übergängen mit
    Exponentialkennlinie den Spannungssprung begrenzen und meldet das; \esim
    erklärt dann keine Konvergenz. Plugins, die das nicht tun, lösen
    unbegrenzt.
\end{description}

Ereignisse (\code{events}) und Sprünge bleiben Sache der Eingabedatei; das
Plugin sollte glatte Kennlinien liefern. Nichtglatte Größen, etwa die Kraft
beim Wechsel einer Netzkante, stören das Newton-Verfahren und die
Schrittweitenregelung; Kennfelder sollten geglättet sein.

\section{Schnappschüsse}
\label{sec:plugin-schnappschuesse}

Ein Plugin mit Orts- oder Gitterdaten (Feldverteilung, Bänder,
Ladungsträgerdichten) kann sie auf Anforderung in Dateien schreiben.

\begin{json}
{ "type": "plugin", "name": "X1", "plugin": "pn", "nodes": ["a", "0"],
  "snapshots": { "times": [1e-6, 5e-6] } }
\end{json}

\begin{itemize}
  \item \code{"snapshots": true}: am Ende jedes Arbeitspunkts und jedes
    Transienten.
  \item \code{{"times": [...]}}: zusätzlich zu diesen Zeiten; \esim legt dort
    einen Zeitpunkt, sodass der Schnappschuss genau dort entsteht.
  \item Die Dateien liegen im Ausgabeordner. Name und Format wählt das Plugin.
    Der Anhang des Namens (\code{<name>.<op | tran.end | tran.t<Zeit>>.<bauteil>})
    kommt von \esim.
  \item In Parameterstudien entstehen keine Schnappschüsse.
\end{itemize}

\section{Sicherheit}
\label{sec:plugin-sicherheit}

Eine Plugin-Bibliothek ist ausführbarer Code mit allen Rechten des Benutzers.
Beim Öffnen einer Eingabedatei mit \code{plugins} oder \code{library} wird die
Bibliothek geladen und ihr Code läuft. Eingabedateien aus unbekannter Quelle
öffnet man deshalb nicht, wenn sie Plugins nennen. \esim schreibt beim Laden in
den Lauf eine Zeile der Form
\code{note: plugin component X1: <Pfad>}; der Editor zeigt sie als Warnung.

\begin{achtung}
  Plugin-Pfade sind Teil der Eingabedatei. Wer eine Datei weitergibt, gibt
  auch den Pfad der Bibliothek weiter; relative Namen vereinfachen das
  Weitergeben. Schalter, die das Laden verhindern (Suchpfad, Betrieb ohne
  Plugins), gibt es noch nicht.
\end{achtung}

\section{Fehlermeldungen}
\label{sec:plugin-fehler}

\begin{longtable}{@{}P{6.2cm}P{7.6cm}@{}}
  \toprule
  Meldung (Auszug) & Ursache \\
  \midrule
  \endhead
  \bottomrule
  \endfoot
  \code{plugin library not found (tried ...)} & Datei fehlt; die Meldung nennt alle versuchten Namen \\
  \code{cannot load plugin '...'} & Das Betriebssystem lädt die Datei nicht: falsche Architektur, fehlende Abhängigkeit \\
  \code{is not an esim plugin} & Die Bibliothek exportiert \code{esim_plugin_entry} nicht \\
  \code{has interface version} & Das Plugin wurde für eine andere Version der Schnittstelle gebaut \\
  \code{plugin "..." is not defined in "plugins"} & Name im Bauteil fehlt in \code{plugins} \\
  \code{input "..."} & Eingang fehlt, ist unbekannt oder sein Ausdruck ist falsch \\
  \code{plugin component 'X1': ...} & Meldung des Plugins, etwa weil die Konfiguration unzulässig ist \\
\end{longtable}

\section{Ein Plugin schreiben}
\label{sec:plugin-schreiben}

Die Schnittstelle ist reines C und steht in einer einzigen Datei,
\datei{plugin/esim_plugin.h}. Weil sie C ist, dürfen Plugin und \esim mit
verschiedenen Compilern gebaut sein, und das Plugin darf in einer anderen
Sprache geschrieben sein. Die Strukturen beginnen mit ihrer Größe, neue
Felder werden nur angehängt; ältere Plugins laufen mit neuerem \esim und
umgekehrt, solange die Versionsnummer \code{api\_version} übereinstimmt.

Ein Plugin exportiert die Funktion \code{esim_plugin_entry()} und liefert
damit eine Tabelle mit

\begin{itemize}
  \item \code{create}, \code{destroy}: Instanz mit der Konfiguration (als
    JSON-Text) anlegen und freigeben,
  \item \code{describe}: Anzahl und Namen der Klemmen, Eingänge, Zweige,
    Ausgänge und inneren Unbekannten, Startwerte, Toleranzen und das feste
    Besetzungsmuster von $\partial f/\partial x$ und $\partial q/\partial x$,
  \item \code{evaluate}: $f$, $q$ und beide Matrizen aus Zeit, Vektor und Modus
    (Gleichstrom, Startwert, Zeitschritt); eine reine Funktion der Eingabe,
  \item optional \code{accept\_step}, \code{next\_breakpoint}, \code{max\_step},
    \code{reset}, \code{last\_error}, \code{noise} und \code{snapshot}.
\end{itemize}

Ein lauffähiges Beispiel mit zwei Modellen (Hubmagnet, Diode mit innerem
Knoten, Begrenzung, Rauschen, Schnappschuss) ist
\datei{plugins/mock/mock_plugin.cpp}; es dient auch als Vorlage. Die
ausführliche Beschreibung der Schnittstelle, ihrer Anforderungen an Feldlöser
und Halbleitersimulatoren und der noch offenen Punkte steht in
\datei{docs/fem-plugin-plan.md} und \datei{docs/halbleiter-plugin.md}.

\begin{hinweis}
  Das Beispielplugin gehört zu den Tests; \esim wird ohne Plugin ausgeliefert.
  Echte Plugins (FEM, Halbleiter) sind eigene Projekte, die nur
  \datei{plugin/esim_plugin.h} brauchen.
\end{hinweis}
